\chapter{The CNJ Model Toolchain}
\label{ch:cnj-model-toolchain}

Direct runtime glTF loading is convenient, but it returns only the first mesh group and fixes
unit scale at~1. The offline \texttt{cna\_tool\_gltf\_to\_cnj} path uses the same semantic
import core and material packing, then writes every group as explicit CNJ and binary sidecars.
It is the reproducible path for multi-skin files, controlled scaling, and inspectable build
artifacts.

\section{A renderer-free build tool}

The tool is configured by \texttt{cmake/ToolGltfToCnj.cmake}, built unconditionally rather than
behind \texttt{CNA\_BUILD\_TESTS}, and links CNA plus sharp-runtime. It does not initialize a
window, graphics device, or renderer. The command line is

\begin{lstlisting}
cna_tool_gltf_to_cnj input.gltf outputDir baseName [unitScale]
\end{lstlisting}

The input may be \texttt{.gltf} or \texttt{.glb}. For each group, the tool writes a model
descriptor and sidecars:

\begin{itemize}
  \item \texttt{<name>.cnj}, a version-2 \texttt{Model} envelope;
  \item \texttt{<name>\_meshN\_verts.bin} and \texttt{\_idx.bin};
  \item optional \texttt{\_morph.bin} files;
  \item optional \texttt{<name>.skeleton.bin};
  \item one version-1 \texttt{AnimationClip} CNJ per clip; and
  \item extracted base, occlusion, and other texture images with stable suffixes.
\end{itemize}

Multiple groups receive stable \texttt{\_static} or skin-name suffixes. Unlike runtime
\texttt{groups.front()}, no later group is silently discarded.

\section{The Model envelope is version 2}

CNJ's common envelope contains a version and a logical type. Most types accept version~1 only;
\texttt{Model} has a per-type maximum of~2. Raising the Model ceiling did not globally bless
version~2 for textures, effects, or animation clips. Validation receives the requested type
and that type's maximum, so a version accepted for one reader can still be rejected by another.

Version~2 adds a top-level \texttt{bones} array. Each record carries a name, parent index, and
16-float local transform, and records are parent-before-child with an identity root at index~0.
Each mesh also receives a \texttt{parentBone} index. These fields preserve the glTF scene
hierarchy and rigid instancing placement that a flat list of pre-transformed meshes could not.

Version-1 Model files remain loadable. When no general hierarchy exists, the reader constructs
the older per-mesh fallback bone. This makes version~2 an additive fidelity improvement rather
than a flag day.

\section{Binary sidecars and bounded reads}

Vertex and index bytes use the packed ABI from Chapter~\ref{ch:gpu-packing}. The descriptor
supplies counts, strides, formats, primitive counts, and paths; the reader verifies limits and
spans before allocating or uploading. Morph sidecars retain base bytes and target deltas used
by \texttt{MorphTargetDataEXT}.

The skeleton sidecar historically contained two matrix blocks: bind poses and inverse binds.
Version~2 can append a third block containing \texttt{SkeletonRootPrefix}. The reader checks
whether at least \texttt{boneCount * 64} bytes remain before reading it; if not, it supplies
identity prefixes. Appending an optional fixed-size block rather than changing the meaning of
the old two is what keeps earlier sidecars valid.

The descriptor may reference separate clip CNJs. Those pass through the manager's normal
cached loading path, so repeated names can share a clip object. Mesh buffers, built-in effects,
and most graph objects are still created per described part and retained by the model's private
ownership bundle.

\section{Offline and runtime parity}

Sharing \texttt{GltfImportCore} reduces drift, but does not prove parity by itself. The two
front ends perform different work after extraction: one creates textures, effects, buffers and
a live model; the other serializes bytes and descriptors that a later CNJ reader recreates.
Tests therefore compare semantic and rendering-visible results across both routes, including
reversed skin order, animation, PBR data, morph tracks, texture extraction, and hierarchy.

The important contract is not byte identity between a live object and JSON. It is that a group
imported directly and the same group converted then loaded from CNJ produce the same scene-node
placement, palette coordinate spaces, vertex/index bytes, material choice, and animation
records. A fix made only in one front end is a regression even if the shared parser still
builds.

\section{Worked example: occlusion through DualTextureEffect}

glTF occlusion uses 1.0 for fully visible light and 0.0 for full occlusion. XNA's
\cnaclass{DualTextureEffect}, however, multiplies its second texture contribution by~2; in that
equation 0.5 is neutral. Passing glTF's occlusion image unchanged would brighten rather than
preserve an unoccluded surface.

For the non-PBR dual-texture approximation, CNA therefore decodes the occlusion image, halves
its RGB channels, and re-encodes it. The resulting mapping is

\[
  o_{dual} = \frac{o_{gltf}}{2},
  \qquad 2o_{dual}=o_{gltf}.
\]

The remapped image is cached separately from the original because the same source image may
also be bound to \cnaclass{PbrEffect}, where glTF's original convention is correct. The runtime
and converter both use this shared extraction rule; tests cover successful remapping and decode
failure.

\section{Choosing a route}

Use direct glTF for rapid iteration on a file known to have one relevant group and unit scale.
Use the converter when the file contains multiple characters or static plus skinned content,
when unit conversion is required, or when asset bytes must be reviewed and committed. Use XNB
when compatibility with an existing XNA content pipeline is the governing requirement. The
legacy \texttt{.skinnedmodel.json} path is for \cnaclass{SkinnedModelEXT}, not an alternate CNJ
spelling.

Do not hand-edit counts or offsets in a generated Model CNJ without regenerating or validating
the sidecars. The JSON is readable, but its numerical metadata and binary files form one asset
contract.
